\chapter*{前言}
\addcontentsline{toc}{chapter}{前言}



写下这份笔记，首先是为了自我提升。SLAM（Simultaneous Localization and Mapping，同时定位与建图）对我而言并不只是一个工程问题：它一方面关心传感器、运动、地图、优化和程序实现，另一方面又不断把人带回一些更基础的问题：什么是状态？什么是观测？什么是约束？为什么位姿不能简单地当作普通欧氏空间中的向量？局部线性化到底是在什么空间里进行的？这些问题看似分散，但背后其实都与集合、映射、拓扑、流形和微分结构有关。

刚开始接触SLAM时，很容易被大量公式、矩阵、Jacobian、李群和优化算法淹没。公式当然重要，但如果只记住计算规则，而没有理解这些规则背后的空间图像，就很难真正知道自己在算什么。比如$\SE(3)$不仅是一组$4\times4$矩阵，它也可以被看作描述机器人位姿的状态空间；扰动量不仅是一个六维向量，它更像是当前位姿附近切空间中的局部坐标；一次优化迭代也不仅是代数更新，而是在非线性空间上借助局部线性结构前进的一小步。写这份笔记，就是希望把这些抽象概念逐渐整理成可以想象、可以追踪、可以和SLAM对象对应起来的数学直观。

这份笔记前半部分建立数学语言与空间直观，后半部分将这些语言依次用于ESKF、IMU预积分、点云配准以及完整SLAM系统。它不试图穷尽每一种实现，而是沿着一条可以逐步追踪的路径展开：
\[
    \begin{aligned}
    \text{逻辑与集合}
    &\longrightarrow\text{拓扑与流形}
    \longrightarrow\text{微分几何与李群}\\
    &\longrightarrow\text{测度与概率}
    \longrightarrow\text{ESKF与相对约束}\longrightarrow\text{SLAM系统}.
    \end{aligned}
\]
这条路径并不是说学习SLAM必须先完整学完这些数学分支，而是说：如果想理解SLAM中那些常见但不太容易说清楚的概念，可以沿着这条线索慢慢建立一套语言。逻辑帮助我们区分命题、定义和推理；集合论让状态、观测、残差、约束这些对象有一个最基本的容器；点集拓扑让我们开始讨论连续性、邻域、紧性和空间结构；拓扑流形告诉我们某些空间整体上可能弯曲复杂，但局部仍然像欧氏空间；可微流形进一步允许我们讨论切空间、微分、Jacobian和局部扰动。最后，这些语言会回到SLAM中的位姿空间、状态估计、误差函数和优化过程。

全书先理解对象，再追踪对象之间的映射，最后推导可计算的公式。集合论部分帮助理解状态集合、观测集合、可行约束和商空间；拓扑与流形部分说明连续性、局部坐标和空间结构；微分几何与李群部分建立切空间、指数映射和局部扰动。第七章沿集合系、可测映射、测度与积分建立概率语言，以少量SLAM例子作对照。第八章用一台带IMU的机器人走完ESKF预测、观测、注入与重置的循环，先说明每个矩阵把什么量送到什么量。第九章再从真实与名义运动出发，推导误差传播、观测更新和重置的关键矩阵。第十章把高频IMU保存为两关键帧之间的预积分约束；第十一章从点云对应与局部协方差出发理解GICP。第十二章把这些估计与约束放回前端、后端和回环的数据流。可以先读第八章建立状态估计的整体认识，再按需要回到第九章追踪计算；读完第十、十一章后，再看第十二章中每份结果怎样被系统使用。

如果把SLAM看成一个工程系统，那么数学提供的并不是唯一的实现方式，而是一种看待问题的结构化视角。它能帮助我们知道哪些量可以相加，哪些量只能在局部相加；哪些误差是在观测空间里度量的，哪些更新是在状态空间附近进行的；为什么同一个机器人姿态可以有不同的坐标表达，为什么优化时常常要在当前估计附近建立局部模型。对我来说，理解这些问题比单纯记住某个公式更重要，因为公式可能会随着模型、传感器和参数化方式变化，而背后的几何直观会一直存在。

阅读这份笔记时，可以不把每个定理都当作需要严格证明的目标，而是把它们看作搭建直观的脚手架。先看清楚对象是什么，再看对象之间如何映射，最后再看这些结构如何落到SLAM问题中。若某些章节暂时显得抽象，也可以先保留印象，等读到位姿空间、切空间或李群时再回头对照。很多数学概念第一次出现时并不自然，但当它们和机器人状态、局部扰动、误差最小化联系起来以后，会变得具体得多。

需要说明的是，写作这份笔记时，笔者仅仅大一，对SLAM和相关数学内容的理解都还很浅。很多内容只是现阶段学习过程中的整理和尝试，并不代表成熟、完整或权威的表述。文中难免存在表述不严谨、理解不充分，甚至理解错误之处。如果读者发现概念混淆、推理漏洞、符号误用或与SLAM实际做法不一致的地方，欢迎及时指正。对我而言，这份笔记本身也是一个持续修正和学习的过程。

最后，希望这份笔记能记录下一个初学者试图把抽象数学和机器人感知问题连接起来的过程。它不追求一步到位，也不追求覆盖所有细节，而是希望在逻辑、集合、拓扑、流形和SLAM之间搭起一条可以继续向前走的路。
